\chapter{Conclusiones}
\label{cap:conclusiones}

El objetivo principal de este trabajo de fin de grado era estudiar si los modelos de lenguaje de gran escala pueden resolver problemas de olimpiadas lingüísticas mediante razonamiento inductivo, sin recibir información explícita previa sobre la lengua analizada. Para ello se ha diseñado, implementado y evaluado un sistema experimental basado en distintas estrategias de \textit{prompting}, aplicado sobre 11 puzzles de traducción del benchmark \textsc{Linguini}, en dirección lengua desconocida $\rightarrow$ inglés, utilizando tres modelos de distinto tamaño: Llama-3.1-8B, Llama-3.3-70B y GPT-OSS-120B. El trabajo parte del baseline directo comparable con el planteamiento de \citet{sanchez2024linguini}, pero incorpora estrategias adicionales orientadas a guiar el razonamiento metalingüístico del modelo mediante diccionarios, reglas revisables, razonamiento humano y autocorrección.

El primer resultado relevante es que la estrategia M2, basada en una presentación incremental acumulativa de los ejemplos, mejora claramente el baseline M1 en el modelo de razonamiento GPT-OSS-120B. En este modelo, M2 obtiene el mejor resultado en 6 de los 11 puzzles evaluados y alcanza una media de chrF++=48.5, frente a 40.6 en M1. Este comportamiento confirma parcialmente la hipótesis central del trabajo: presentar los ejemplos de forma progresiva y obligar al modelo a revisar su diccionario y sus reglas puede mejorar el razonamiento inductivo frente a una presentación completa de una sola vez. Sin embargo, el beneficio no es uniforme. En Llama-3.1-8B, M2 baja ligeramente respecto a M1, y en Llama-3.3-70B también queda por debajo del baseline. Esto indica que la estrategia incremental no es universalmente mejor, sino que requiere que el modelo pueda mantener hipótesis acumuladas, revisar contradicciones y no degradar el conocimiento construido en pasos anteriores.

El segundo hallazgo, quizá el más sorprendente, es que el modelo con mejor rendimiento global no es el de mayor tamaño. Llama-3.3-70B obtiene la mejor media de resultados máximos por puzzle, con chrF++=55.9, frente a 52.3 de GPT-OSS-120B y 43.4 de Llama-3.1-8B. Además, consigue el mejor resultado en 7 de los 11 puzzles. Este resultado invierte la jerarquía esperada por tamaño y sugiere que la capacidad de razonamiento inductivo lingüístico no escala de forma lineal con el número de parámetros, al menos dentro de los modelos y proveedores evaluados. Otros factores como el ajuste instructivo, el formato de salida, la estabilidad de generación o la gestión del contexto parecen tener un peso importante, al menos en las pruebas realizadas hasta ahora.

El tercer resultado importante se observa en M5, la estrategia de autocorrección iterativa. Aunque la hipótesis inicial era que una segunda fase de revisión podría corregir contradicciones internas, los resultados muestran que la autocorrección sin retroalimentación externa tiende a empeorar el rendimiento en la mayoría de casos. El efecto es especialmente claro en Llama-3.3-70B, donde algunos primeros borradores muy buenos se degradan tras la revisión: Hankun baja de 77.8 a 26.6, Lakota de 69.2 a 28.0 y Nubio kunuz de 47.7 a 16.6. Este comportamiento es coherente con las limitaciones del paradigma \textit{Self-Refine} descritas por \citet{madaan2023selfrefine}: si la crítica procede del propio modelo y no existe una señal externa de validación, la revisión puede reforzar errores iniciales en lugar de corregirlos. Aquí surge el efecto contrario, el problema se agrava porque el modelo puede usar sus propias reglas como árbitro, incluso cuando esas reglas no están suficientemente justificadas por los ejemplos de entrenamiento.

Los experimentos complementarios permiten matizar todavía más la respuesta a la pregunta de investigación. La dirección inversa, inglés $\rightarrow$ lengua desconocida, resulta considerablemente más difícil que la dirección principal lengua desconocida $\rightarrow$ inglés. Además, las tareas morfológicas y numéricas refuerzan esta conclusión: en Guazacapán Xinka todos los modelos obtienen solo 1/4 en Exact Match, y en Supyire ningún modelo alcanza Exact Match en las tareas numéricas. Los modelos identifican regularidades parciales, pero no siempre logran aplicarlas con la precisión formal necesaria.

En conjunto, los resultados muestran que los modelos de lenguaje sí poseen cierta capacidad metalingüística inducida por contexto: pueden extraer vocabulario, detectar regularidades morfológicas, formular hipótesis y aplicar algunas reglas a frases nuevas. No obstante, esa capacidad no alcanza el nivel humano de forma estable en ninguna configuración evaluada. Las estrategias de prompting influyen de manera significativa en el rendimiento: según el modelo y el puzzle, cambiar la metodología puede suponer diferencias de más de 30 puntos. 
Aunque los modelos actuales todavía presentan limitaciones importantes, especialmente en producción morfológica y revisión autónoma, los resultados sugieren que este tipo de enfoques sí puede ser útil para analizar lenguas de baja densidad de recursos en contextos donde solo se dispone de unos pocos ejemplos.